\section{LENS}

\subsection{Design Goal}
LENS is designed to restore metric separability in multi-subject personalized generation. It is not sufficient for a metric to output lower scores as subject count increases. A useful evaluator must still distinguish obviously better generations from obviously worse ones under the same prompt-reference condition, especially at high subject counts where conventional metrics collapse.

\subsection{Model Interface}
LENS takes as input the reference subjects, the prompt, and a generated image pair. The references define the identity targets; the prompt defines the semantic and interaction constraints; the two generated images provide the paired comparison for ranking. The output contains:
\begin{itemize}
    \item a scalar preference score for each image,
    \item a binary diagnostic vector for each image over Existence, Appearance, and Interaction.
\end{itemize}

\subsection{Architecture}
The evaluator uses a multimodal backbone and attaches two lightweight heads on top of the shared representation:
\begin{itemize}
    \item a \textbf{score head} that outputs a scalar score,
    \item a \textbf{classification head} that outputs the three diagnostic logits.
\end{itemize}

This dual-head design is important. The score head captures fine-grained ordering information that binary class labels cannot represent, while the classification head regularizes the shared representation toward interpretable failure modes.

\subsection{Why the Binary Taxonomy is Enough}
At first glance, a binary diagnostic scheme may appear too coarse for nuanced quality assessment. However, LENS does not rely on the diagnostic head alone. The role of the diagnostic head is to define \emph{what kind of error is present}. The role of the score head is to define \emph{how bad the image is relative to another image}. This separation allows the training objective to remain simple for annotators while preserving continuous score discrimination in the learned metric.

\subsection{Expected Failure Sensitivity}
The main expected advantage of LENS over global metrics is its ability to remain reference-aware and locally diagnostic when $N$ becomes large. If one generator preserves most identities while another collapses several of them, LENS should assign a higher preference score to the first image even when conventional metrics report nearly overlapping values. This property is the core criterion for success in the final evaluation.
